\documentclass[11pt]{article}
\usepackage{amsmath,amssymb,hyperref}
\title{A retirement proof does not have to arrive after the fence}
\date{}
\begin{document}
\maketitle

\paragraph{Cut and service.}
An incoming controller must perform one new effect on a shared fixture after the old scope has irrevocably stopped issuing and its effects using that fixture have settled. Its own independent staging work takes $p$ unit ticks starting at cut zero. The new effect then takes one tick. An available staging site and permission to use it are prerequisites, not furnished by the proof. An old effect away from the fixture must still count if it can later arrive there. No competing scope can acquire a new permit under this envelope.

Old issues have an immutable numbering from 1. The effect-side issuer supplies a verified fence carrying scope $s$ (epoch, issuer and fixture) and final issued number $N$. This proves both that all outstanding old operations are represented and that the old scope cannot issue again. An effect-side certificate $(s',C)$ attests actual terminal settlement of every issued operation through $C$ for scope $s'$. Terminal settlement is permanent and includes cancellation only after actual drain. A delivery or start acknowledgment does not attest this. A stalled operation cannot be skipped by reporting a larger individually settled ID as a prefix. Number reuse, incompletely fenced writers, and unrecorded effects are outside this model.

\paragraph{Join rather than serialize.}
Upon receipt of the final fence, a valid clearance proves the old obligation discharged if $s=s'$ and $C\ge N$. Every old operation is enumerated and closed. The requirement is necessary as an evidence law for the modeled opaque suffix: if $C<N$, an unresolved issued operation at $N$ is not excluded by the proof. Another source with trustworthy time bounds, complete additional observations, or no capacity for such an issue may admit another proof. Missing the fence leaves a possible later old issue even after an empty reading.

A clearance can be sampled and delivered before the fence. Only coverage of the final number for the same scope matters. A reading covering the current issue frontier does not authorize entry before the actual fence. This is a local join of verified evidence, not reconstruction of arbitrary global histories. Equal numbers, matching timestamps, or empty software queues do not supply the scope or physical interpretation.

Let $g$ be the first cut at which the final fence and a covering clearance are both recognized. With independent preparation the sharp finish is $\max(p,g)+1$: do preparation from zero and apply the joined proof at entry. The evidence and physical laws give the two necessary lower bounds. Starting preparation only after $g$ instead takes $g+p+1$, a loss of $\min(p,g)$. If preparation consumes the held fixture, that overlap is unavailable. Demanding a new sample made after recognition of the fence can also delay entry despite already-held valid evidence. These are separate choices; the theorem concerns one effect on one fixture and does not forecast a new obligation.

\paragraph{Four source controls.}
With $p=2$, one old operation is settled and reported at tick 1. The fence is recognized at tick 3 and proves final number 1. The new effect finishes at tick 4. Requiring a later sample (at tick 5) or serial preparation each takes six ticks. With the same early clearance but another old issue, the fence instead names number 2. That extra issue physically settles at 5 and is reported at 6; the correct new finish is 7. An early report with the right count from a different scope also cannot replace the late proof. With the fence report erased, physical emptiness does not provide the required evidence. Discarding old issued effects from the controller state can therefore approve an entry the plant cannot safely execute.

All 1,392 feasibility queries matched a raw decision-tree solver and the pinned upstream safety fixed point and action-rectangle methods from \href{https://github.com/nightly/contract-shielding/blob/07ef8d482325ce8450753cba0276511b81d26ec3/src/shield/_assume_guarantee.py}{Adalat et al.'s implementation}. The adapter uses the live source transcript, includes missed deadlines as bad states, and gives the old side its prescribed physical trajectory. This tests the same-information product kernel, not the complete trained model or a two-controller synthesis. The paper's \href{https://arxiv.org/html/2606.14130v2}{Theorem 5.2} concerns selection between covered episode starts; its product-graph implementation already accepts caller-supplied product state. Overlap beats serialized preparation, not the correctly seeded same-information implementation. Lean checks scoped coverage, the exact opaque-prefix bound, the timing difference, and a minimal counterexample.

\paragraph{Relation to existing results.}
\href{https://dspace.mit.edu/entities/publication/80846d94-fcd3-40e6-87fb-8d91fe99a5d1}{Liskov and Cowling's viewstamped reconfiguration} is a precedent for fencing an old sequence and carrying it to a new configuration. \href{https://www.rfc-editor.org/rfc/rfc9293.html}{TCP} cumulatively acknowledges bytes received, not independent physical effects. The local proof consumes the contiguous actual settlement evidence defined in the prior \href{https://github.com/eternal-november/eos-research-programmes/pull/359}{swarm receipt result}; it is not a claim of novelty for sequence-number fences. In the pinned \href{https://github.com/eternal-november/eos-research-programmes/pull/418}{RWARE trace}, a delivered shelf remains attached to the old vehicle at the physical outlet even after the queue advances. That is a negative source control, not a native fence. The positive model is a bounded R2/R5 comparison, not independent integration or full acceptance.
\end{document}
